% --- CORRECCIÓN IMPORTANTE EN LA PRIMERA LÍNEA ---
% Agregamos "xcolor=table" para que funcione \rowcolor sin errores
\documentclass[aspectratio=169, 10pt, xcolor=table]{beamer}

% --- Configuración del Tema (estilo profesional como el ejemplo) ---
\usetheme{Madrid}
\usecolortheme{dolphin}

% --- Paquetes necesarios ---
\usepackage[utf8]{inputenc}
\usepackage[spanish]{babel}
\usepackage{amsmath, amssymb}
\usepackage{booktabs}
\usepackage{graphicx}
\usepackage{listings}
\usepackage{tikz}
\usetikzlibrary{arrows.meta, positioning, shapes.geometric}

% --- Ruta de Imágenes (crea la carpeta "imagenes" y coloca ahí tus archivos) ---
\graphicspath{{./imagenes/}}

% --- Estilo para código Python ---
\definecolor{codegreen}{rgb}{0,0.6,0}
\definecolor{codegray}{rgb}{0.5,0.5,0.5}
\definecolor{codepurple}{rgb}{0.58,0,0.82}
\definecolor{backcolour}{rgb}{0.95,0.95,0.92}
\lstset{
    language=Python,
    backgroundcolor=\color{backcolour},
    commentstyle=\color{codegreen},
    keywordstyle=\color{blue},
    numberstyle=\tiny\color{codegray},
    stringstyle=\color{codepurple},
    basicstyle=\ttfamily\scriptsize,
    breaklines=true,
    frame=single,
    showstringspaces=false
}

% --- Pie de página personalizado ---
\makeatletter

\setbeamertemplate{footline}{
  \leavevmode%
  \hbox{%
  \begin{beamercolorbox}[wd=.333333\paperwidth,ht=2.25ex,dp=1ex,center]{author in head/foot}%
    \usebeamerfont{author in head/foot}\insertshortauthor
  \end{beamercolorbox}%
  \begin{beamercolorbox}[wd=.333333\paperwidth,ht=2.25ex,dp=1ex,center]{title in head/foot}%
    \usebeamerfont{title in head/foot}Sesión 9
  \end{beamercolorbox}%
  \begin{beamercolorbox}[wd=.333333\paperwidth,ht=2.25ex,dp=1ex,right]{date in head/foot}%
    \usebeamerfont{date in head/foot}Machine Learning \hfill \insertframenumber/\inserttotalframenumber \hspace*{0.3cm}
  \end{beamercolorbox}}%
  \vskip0pt%
}

\makeatother

% --- Datos de la portada ---
\title[SESIÓN 9]{\textbf{Sesión 9}}
\subtitle{Aprendizaje No Supervisado: \\ Clustering y Análisis de Componentes Principales}
\author[Prof. C. Sánchez]{Carlos César Sánchez Coronel}
\institute{\textbf{MACHINE LEARNING}}
\date{}

% Logo opcional (descomentar si tienes la imagen)
% \titlegraphic{\includegraphics[height=1.5cm]{logo.png}}

% --- Eliminar la barra de navegación (opcional) ---
\setbeamertemplate{navigation symbols}{}

% --- Índice al inicio de cada sección ---
\AtBeginSection[]{
  \begin{frame}
    \frametitle{Contenido}
    \tableofcontents[currentsection]
  \end{frame}
}

% ============================================================================
% INICIO DEL DOCUMENTO
% ============================================================================
\begin{document}

% --- Portada ---
\begin{frame}
    \titlepage
\end{frame}

% --- Índice general ---
\begin{frame}{Contenido}
    \tableofcontents
\end{frame}

% ============================================================================
% SECCIÓN 1: INTRODUCCIÓN
% ============================================================================
\section{Introducción al Aprendizaje No Supervisado}

\begin{frame}{¿Qué es el aprendizaje no supervisado?}
    \begin{itemize}
        \item Trabajamos con datos \textbf{sin etiquetas}.
        \item Objetivos principales:
        \begin{itemize}
            \item Descubrir grupos o segmentos naturales (\textbf{clustering}).
            \item Reducir la dimensionalidad para visualización, compresión o preprocesamiento (\textbf{PCA}).
            \item Detectar anomalías o valores atípicos.
        \end{itemize}
        \item No hay una métrica de rendimiento única; la evaluación depende del objetivo y del conocimiento del dominio.
    \end{itemize}
    % Basado en introducción de la teoría
\end{frame}

% ============================================================================
% SECCIÓN 2: ANÁLISIS DE COMPONENTES PRINCIPALES (PCA)
% ============================================================================
\section{Análisis de Componentes Principales (PCA)}

\begin{frame}{PCA: fundamento matemático}
    \begin{block}{Objetivo}
        Proyectar los datos en un espacio de menor dimensión maximizando la varianza retenida.
    \end{block}
    \begin{enumerate}
        \item Centrar la matriz de datos $\mathbf{X}$ (restar media por columna).
        \item Calcular matriz de covarianza: $\displaystyle \mathbf{\Sigma} = \frac{1}{n-1}\mathbf{X}^T\mathbf{X}$.
        \item Obtener autovalores $\lambda_1 \ge \lambda_2 \ge \dots \ge \lambda_p$ y autovectores asociados $\mathbf{v}_1, \dots, \mathbf{v}_p$.
        \item El $k$-ésimo componente principal es $\mathbf{X}\mathbf{v}_k$.
    \end{enumerate}
    % Basado en Pregunta 1
    % IMAGEN: Esquema de proyección PCA
\end{frame}

\begin{frame}{Interpretación de componentes}
    \begin{itemize}
        \item \textbf{Autovalores}: representan la varianza explicada por cada componente.
        \item \textbf{Proporción de varianza explicada}: $\displaystyle \frac{\lambda_j}{\sum_{i=1}^p \lambda_i}$.
        \item \textbf{Loadings} (autovectores): indican la contribución de cada variable original al componente.
        \item Un valor alto en magnitud señala variables importantes para ese componente.
    \end{itemize}
    % Basado en Pregunta 2
\end{frame}

\begin{frame}{Scree plot y elección del número de componentes}
    \begin{columns}[t]
        \column{0.5\textwidth}
        \textbf{Scree plot}: gráfico de autovalores en orden descendente.
        \begin{itemize}
            \item Se busca un ``codo'' donde la caída se vuelve marginal.
            \item Se retienen los componentes hasta el codo.
        \end{itemize}
        \column{0.5\textwidth}
        % IMAGEN: Ejemplo de scree plot
    \end{columns}
    \vspace{0.3cm}
    \begin{exampleblock}{Ejemplo numérico}
        Autovalores: [4.5, 2.1, 1.3, 0.8, 0.6, 0.4, 0.3, 0.2, 0.1, 0.05] \\
        Varianza total = 10.35. \\
        Primer componente: 43.5\%, dos primeros: 63.8\%. El codo sugiere k=3.
    \end{exampleblock}
    % Basado en Pregunta 2
\end{frame}

\begin{frame}{Aplicaciones de PCA}
    \begin{itemize}
        \item \textbf{Visualización}: proyectar datos a 2D o 3D para inspección.
        \item \textbf{Reducción de ruido}: descartar componentes de baja varianza (ruido).
        \item \textbf{Preprocesamiento}: reducir dimensionalidad antes de modelos sensibles (KNN, clustering).
        \item \textbf{Compresión}: almacenar solo las primeras componentes.
    \end{itemize}
    % Basado en Pregunta 3
\end{frame}

% ============================================================================
% SECCIÓN 3: K-MEANS CLUSTERING
% ============================================================================
\section{K-Means Clustering}

\begin{frame}{Algoritmo de K-Means}
    \begin{enumerate}
        \item Inicializar $k$ centroides (aleatorio o K-Means++).
        \item Asignar cada punto al centroide más cercano (distancia euclidiana).
        \item Recalcular cada centroide como la media de los puntos asignados.
        \item Repetir 2 y 3 hasta convergencia (centroides no cambian o máximo de iteraciones).
    \end{enumerate}
    \vspace{0.3cm}
    \begin{block}{Función objetivo (inercia)}
        \[
        J = \sum_{j=1}^{k} \sum_{x \in C_j} \|x - \mu_j\|^2
        \]
        Minimiza la suma de distancias cuadráticas intra-cluster.
    \end{block}
    % Basado en Pregunta 4
    % IMAGEN: Ilustración de iteraciones de K-Means
\end{frame}

\begin{frame}{Elección del número de clusters $k$}
    \begin{columns}[t]
        \column{0.5\textwidth}
        \textbf{Método del codo}
        \begin{itemize}
            \item Graficar inercia vs $k$.
            \item Elegir $k$ donde la mejora se vuelve marginal.
        \end{itemize}
        \column{0.5\textwidth}
        \textbf{Coeficiente de silueta}
        \begin{itemize}
            \item Mide cohesión y separación.
            \item Promedio para cada $k$; se elige el $k$ con mayor valor.
        \end{itemize}
    \end{columns}
    \vspace{0.3cm}
    \begin{exampleblock}{Ejemplo (Pregunta 5)}
        Inercia: [1000,600,450,400,380,370,365,360,358] \\
        Silueta: [0.35,0.48,0.52,0.50,0.47,0.44,0.41,0.38,0.35] \\
        Codo sugiere $k=3$ o $4$; silueta máxima en $k=3$. Se elige $k=3$.
    \end{exampleblock}
\end{frame}

\begin{frame}{Inicialización: K-Means++}
    \begin{itemize}
        \item K-Means es sensible a la inicialización (óptimos locales).
        \item \textbf{K-Means++} elige centroides dispersos:
        \begin{enumerate}
            \item Primer centroide aleatorio.
            \item Para cada punto, calcular distancia al centroide más cercano ya elegido.
            \item Elegir siguiente centroide con probabilidad proporcional al cuadrado de esa distancia.
            \item Repetir hasta tener $k$ centroides.
        \end{enumerate}
        \item Es la implementación por defecto en scikit-learn.
    \end{enumerate}
    % Basado en Pregunta 4
\end{frame}

\begin{frame}{Limitaciones de K-Means}
    \begin{itemize}
        \item Asume clusters \textbf{esféricos} y de tamaño similar.
        \item Sensible a \textbf{outliers}.
        \item Requiere especificar $k$.
        \item Puede converger a \textbf{óptimos locales}.
    \end{itemize}
    \vspace{0.3cm}
    \begin{exampleblock}{Aplicaciones típicas}
        Segmentación de clientes, compresión de imágenes, agrupamiento de documentos.
    \end{exampleblock}
    % Basado en Pregunta 6
\end{frame}

% ============================================================================
% SECCIÓN 4: DBSCAN
% ============================================================================
\section{DBSCAN (Density-Based Clustering)}

\begin{frame}{DBSCAN: conceptos básicos}
    \begin{block}{Parámetros}
        \begin{itemize}
            \item $\varepsilon$ (eps): radio de vecindad.
            \item \texttt{minPts}: número mínimo de puntos para formar una región densa.
        \end{itemize}
    \end{block}
    \begin{block}{Clasificación de puntos}
        \begin{itemize}
            \item \textbf{Punto núcleo}: tiene al menos \texttt{minPts} vecinos dentro de $\varepsilon$.
            \item \textbf{Punto borde}: no es núcleo pero es vecino de un núcleo.
            \item \textbf{Ruido}: no es núcleo ni borde.
        \end{itemize}
    \end{block}
    % Basado en Pregunta 7
    % IMAGEN: Diagrama de puntos núcleo, borde y ruido
\end{frame}

\begin{frame}{Ventajas y desventajas de DBSCAN}
    \begin{columns}[t]
        \column{0.5\textwidth}
        \textbf{Ventajas}
        \begin{itemize}
            \item No requiere $k$ previo.
            \item Detecta clusters de forma arbitraria.
            \item Identifica outliers como ruido.
        \end{itemize}
        \column{0.5\textwidth}
        \textbf{Desventajas}
        \begin{itemize}
            \item Sensible a $\varepsilon$ y \texttt{minPts}.
            \item Difícil de ajustar si las densidades varían.
        \end{itemize}
    \end{columns}
    \vspace{0.3cm}
    \begin{exampleblock}{Aplicaciones}
        Detección de anomalías, datos geoespaciales, formas no convexas.
    \end{exampleblock}
    % Basado en Pregunta 8
\end{frame}

% ============================================================================
% SECCIÓN 5: CLUSTERING JERÁRQUICO
% ============================================================================
\section{Clustering Jerárquico}

\begin{frame}{Clustering jerárquico aglomerativo}
    \begin{enumerate}
        \item Comenzar con cada punto como un cluster individual.
        \item En cada paso, fusionar los dos clusters más cercanos según una medida de \textbf{enlace}.
        \item Repetir hasta un solo cluster.
    \end{enumerate}
    \vspace{0.3cm}
    \begin{block}{Tipos de enlace}
        \begin{itemize}
            \item \textbf{Single linkage}: distancia mínima entre puntos → forma cadenas.
            \item \textbf{Complete linkage}: distancia máxima → clusters compactos.
            \item \textbf{Average linkage}: distancia promedio.
            \item \textbf{Ward}: minimiza el incremento de varianza intra-cluster.
        \end{itemize}
    \end{block}
    % Basado en Pregunta 9
\end{frame}

\begin{frame}{Dendrograma y elección de clusters}
    \begin{columns}[t]
        \column{0.5\textwidth}
        \textbf{Dendrograma}
        \begin{itemize}
            \item Representa la jerarquía de fusiones.
            \item Cortando a cierta altura se obtienen clusters.
            \item La altura del corte se elige según la distancia de fusión (buscar un salto grande).
        \end{itemize}
        \column{0.5\textwidth}
        % IMAGEN: Dendrograma de ejemplo
    \end{columns}
    \vspace{0.3cm}
    \begin{itemize}
        \item Ventaja: no requiere $k$ previo, jerarquía interpretable.
        \item Desventaja: costoso computacionalmente ($O(n^3)$ o $O(n^2)$).
    \end{itemize}
    % Basado en Pregunta 9
\end{frame}

% ============================================================================
% SECCIÓN 6: MÉTRICAS DE EVALUACIÓN PARA CLUSTERING
% ============================================================================
\section{Métricas de Evaluación para Clustering}

\begin{frame}{Coeficiente de silueta}
    Para cada punto $i$:
    \begin{itemize}
        \item $a(i)$: distancia media a los puntos de su mismo cluster (cohesión).
        \item $b(i)$: distancia media al cluster vecino más cercano (separación).
    \end{itemize}
    \[
    s(i) = \frac{b(i) - a(i)}{\max\{a(i), b(i)\}}
    \]
    \begin{itemize}
        \item $s(i) \approx 1$: cluster compacto y bien separado.
        \item $s(i) \approx 0$: solapamiento.
        \item $s(i) < 0$: asignación incorrecta.
    \end{itemize}
    Se promedia sobre todos los puntos.
    % Basado en Pregunta 10
\end{frame}

\begin{frame}{Índice de Davies-Bouldin}
    \begin{itemize}
        \item Mide la similitud promedio entre cada cluster y su más similar.
        \item Para cada cluster $i$, se define $R_i = \max_{j \neq i} \frac{s_i + s_j}{d_{ij}}$, donde $s_i$ es la dispersión intra-cluster (distancia media al centroide) y $d_{ij}$ es la distancia entre centroides.
        \item Índice final = $\frac{1}{k}\sum_{i=1}^k R_i$.
        \item \textbf{Valores más bajos indican mejor separación}.
    \end{itemize}
    % Basado en Pregunta 10
\end{frame}

\begin{frame}{Comparación de algoritmos de clustering}
    \begin{table}
        \centering
        \begin{tabular}{l|c|c|c}
            \toprule
            \textbf{Algoritmo} & \textbf{Forma clusters} & \textbf{Número de clusters} & \textbf{Ruido} \\
            \midrule
            K-Means & Esféricos & Fijo $k$ & No \\
            DBSCAN & Arbitraria & Automático & Sí \\
            Jerárquico & Cualquier & Corte en dendrograma & No \\
            \bottomrule
        \end{tabular}
    \end{table}
    % Basado en teoría
\end{frame}

% ============================================================================
% SECCIÓN 7: CASO INTEGRADO: SEGMENTACIÓN DE CLIENTES
% ============================================================================
\section{Caso Integrado: Segmentación de Clientes}

\begin{frame}{Problema y datos}
    \begin{itemize}
        \item Empresa de retail con 5000 clientes.
        \item Variables: monto total gastado, frecuencia de compra, antigüedad, edad.
        \item Objetivo: segmentar para personalizar campañas de marketing.
    \end{itemize}
    % Basado en Pregunta 15
\end{frame}

\begin{frame}{Pipeline de segmentación}
    \begin{enumerate}
        \item \textbf{Preprocesamiento}: estandarización (Z-score) de variables numéricas.
        \item \textbf{PCA}: reducir dimensionalidad para visualización (elegir componentes que expliquen al menos 75\% de varianza).
        \item \textbf{K-Means}: probar $k$ de 2 a 10, usar método del codo y silueta.
        \item \textbf{Seleccionar $k$ óptimo} (ej. $k=4$).
        \item \textbf{Caracterizar clusters} mediante perfiles medios en variables originales.
    \end{enumerate}
\end{frame}

\begin{frame}{Resultados y acciones de marketing}
    \begin{table}
        \centering
        \begin{tabular}{l|c|c|c}
            \toprule
            \textbf{Cluster} & \textbf{Edad} & \textbf{Ingresos} & \textbf{Frecuencia} \\
            \midrule
            0 & Baja & Altos & Alta \\
            1 & Media-alta & Bajos & Baja \\
            2 & Alta & Medios & Baja \\
            3 & Baja & Bajos & Media \\
            \bottomrule
        \end{tabular}
    \end{table}
    \begin{itemize}
        \item Cluster 0: ofertas exclusivas, productos premium.
        \item Cluster 1: promociones para incentivar compras.
        \item Cluster 2: productos adaptados a mayores.
        \item Cluster 3: descuentos para jóvenes.
    \end{itemize}
    % Basado en Pregunta 14
\end{frame}

% ============================================================================
% SECCIÓN 8: IMPLEMENTACIÓN EN PYTHON
% ============================================================================
\section{Implementación en Python}

\begin{frame}[fragile]{PCA con scikit-learn}
    \begin{lstlisting}
from sklearn.decomposition import PCA
from sklearn.preprocessing import StandardScaler

scaler = StandardScaler()
X_scaled = scaler.fit_transform(X)

pca = PCA(n_components=2)
X_pca = pca.fit_transform(X_scaled)

print("Varianza explicada:", pca.explained_variance_ratio_)
    \end{lstlisting}
\end{frame}

\begin{frame}[fragile]{K-Means y selección de $k$}
    \begin{lstlisting}
from sklearn.cluster import KMeans
from sklearn.metrics import silhouette_score

inertias, silhouettes = [], []
for k in range(2, 11):
    kmeans = KMeans(n_clusters=k, random_state=42)
    labels = kmeans.fit_predict(X_scaled)
    inertias.append(kmeans.inertia_)
    silhouettes.append(silhouette_score(X_scaled, labels))

# Elegir k con mayor silueta
best_k = np.argmax(silhouettes) + 2
kmeans = KMeans(n_clusters=best_k, random_state=42)
labels = kmeans.fit_predict(X_scaled)
    \end{lstlisting}
\end{frame}

\begin{frame}[fragile]{DBSCAN y clustering jerárquico}
    \begin{lstlisting}
from sklearn.cluster import DBSCAN

dbscan = DBSCAN(eps=0.5, min_samples=5)
labels_db = dbscan.fit_predict(X_scaled)  # -1 indica ruido

# Clustering jerárquico y dendrograma
from scipy.cluster.hierarchy import dendrogram, linkage
import matplotlib.pyplot as plt

Z = linkage(X_scaled, method='ward')
plt.figure(figsize=(10,5))
dendrogram(Z)
plt.show()
    \end{lstlisting}
\end{frame}

% ============================================================================
% SECCIÓN 9: CONEXIÓN Y RESUMEN
% ============================================================================
\section{Conexión y Resumen}

\begin{frame}{Conexión con el resto del curso}
    \begin{itemize}
        \item PCA se usa como preprocesamiento en modelos supervisados para reducir dimensionalidad y ruido.
        \item Los clusters pueden servir como nuevas características en modelos supervisados.
        \item DBSCAN es base para detectores de anomalías.
        \item La interpretación de clusters complementa el análisis exploratorio.
    \end{itemize}
\end{frame}

\begin{frame}{Lo que hemos aprendido}
    \begin{exampleblock}{Conceptos clave}
        \begin{itemize}
            \item \textbf{PCA}: reducción de dimensionalidad, autovalores, scree plot.
            \item \textbf{K-Means}: algoritmo, inercia, método del codo, silueta, K-Means++.
            \item \textbf{DBSCAN}: densidad, puntos núcleo/borde/ruido, parámetros.
            \item \textbf{Clustering jerárquico}: aglomerativo, enlaces, dendrograma.
            \item \textbf{Métricas}: silueta, Davies-Bouldin.
            \item \textbf{Caso práctico}: segmentación de clientes.
        \end{itemize}
    \end{exampleblock}
\end{frame}

% --- Diapositiva final ---
\begin{frame}
    \centering
    \Huge \textbf{¡Gracias!}

\end{frame}

\end{document}